\section{Additional Results}
\label{sec:additional-results}

% 附录 B：扩展实验结果（参考 SWE-bench Appendix C "Additional Results"）
% 附录 B：扩展实验结果（参考 SWE-bench Appendix C "Additional Results"）

\subsection{Full Per-Domain Results}
\label{sec:full-per-domain}
% \subsection{完整的各领域结果}

Table~\ref{tab:full-domain-results} reports the full per-domain breakdown from the 2026-11-09 frozen evaluation (\Neval{} runs across all 13 domains, DeepSeek-V4-Pro, Cline v3.0.51, \texttt{thinking=low}, 8 workers). The overall Pass@1 is 10.0\% (62 PASS / 620 decided runs); 70 ERROR and 6 ABSENT runs are excluded from the denominator but reported per domain. Column sums are 62/558/70/6/696, and for every row PASS$+$FAIL$+$ERROR$+$ABSENT$=$Total.
% 表~\ref{tab:full-domain-results} 报告 2026-11-09 冻结评测（\Neval{} 个 run，覆盖全部 13 个领域，
% DeepSeek-V4-Pro、Cline v3.0.51、\texttt{thinking=low}、8 workers）的完整各领域分解。
% 总体 Pass@1 为 10.0\%（62 PASS / 620 已判定 run）；70 个 ERROR 与 6 个 ABSENT 不计入分母但逐域列出。
% 列合计 62/558/70/6/696，且每行 PASS+FAIL+ERROR+ABSENT=Total。

\begin{table}[ht]
\centering
\small
\caption{Full Pass@1 (\%) by domain on the 2026-11-09 full evaluation (696 cases, DeepSeek-V4-Pro, Cline v3.0.51, \texttt{thinking=low}, 8 workers, 2400s timeout). Pass\% = PASS/(PASS+FAIL); ERROR and ABSENT are excluded from the denominator.}
% \caption{2026-11-09 完整评测（696 例，DeepSeek-V4-Pro、Cline v3.0.51、\texttt{thinking=low}、8 workers、超时 2400s）各领域 Pass@1 (\%)。Pass\% = PASS/(PASS+FAIL)；ERROR 与 ABSENT 不计入分母。}
\label{tab:full-domain-results}
\scriptsize
\begin{tabular}{@{}lcccccc@{}}
\toprule
\textbf{Domain} & \textbf{PASS} & \textbf{FAIL} & \textbf{ERROR} & \textbf{ABSENT} & \textbf{Total} & \textbf{Pass\%} \\
\midrule
schematic\_netlist & 13 & 100 & 18 & 0 & 131 & 11.5 \\
spice\_sim         & 8  & 87  & 7  & 2 & 104 & 8.4  \\
physical\_verify   & 2  & 66  & 4  & 1 & 73  & 2.9  \\
digital            & 5  & 40  & 6  & 0 & 51  & 11.1 \\
analog\_context    & 4  & 39  & 12 & 0 & 55  & 9.3  \\
waveform           & 2  & 37  & 8  & 0 & 47  & 5.1  \\
layout\_virtuoso   & 4  & 35  & 0  & 1 & 40  & 10.3 \\
layout\_aether     & 1  & 33  & 4  & 0 & 38  & 2.9  \\
verification       & 3  & 30  & 5  & 0 & 38  & 9.1  \\
eda\_other         & 6  & 27  & 4  & 1 & 38  & 18.2 \\
pg\_layout         & 4  & 24  & 0  & 0 & 28  & 14.3 \\
characterization   & 5  & 21  & 2  & 0 & 28  & 19.2 \\
pnr\_signoff       & 5  & 19  & 0  & 1 & 25  & 20.8 \\
\midrule
\textbf{Total}     & \textbf{62} & \textbf{558} & \textbf{70} & \textbf{6} & \textbf{696} & \textbf{10.0} \\
\bottomrule
\end{tabular}
\end{table}

\subsection{Per-Model Results}
\label{sec:per-model}
% \subsection{各模型结果}

This submission reports a \emph{single-backbone} evaluation: only DeepSeek-V4-Pro
has been run to completion over the full \Neval{}-run snapshot (10.0\% Pass@1,
62/620 decided). GLM-5.2, Kimi-K2.6, DeepSeek-V4-Flash and Doubao-Seed-2.0-Pro
are configured in the same harness but their sweeps are still running, so we
report no numbers for them rather than a table of dashes. At the observed batch
cost (696 runs in $\approx$24\,h at 8 workers), the remaining four backbones
amount to roughly five machine-days; we will report the sweep in the
camera-ready version and repeat it as the multi-backbone test of RQ3
(Section~\ref{sec:model-scaling}).
% 本投稿给出的是\emph{单骨干}评测：只有 DeepSeek-V4-Pro 在完整的 \Neval{}-run 快照上跑到完成
% （10.0\% Pass@1，62/620 已判定）。GLM-5.2、Kimi-K2.6、DeepSeek-V4-Flash、Doubao-Seed-2.0-Pro
% 已接入同一框架但扫描仍在跑，因此这里不报数字，也不放一张全是 "---" 的表。
% 按本批实测成本（696 run / 8 workers ≈ 24 小时），余下四个骨干约需五个机器日；camera-ready
% 版本将给出该扫描，并用它作为 RQ3 的多骨干检验。

\subsection{Failure Taxonomy by Domain}
\label{sec:failure-by-domain}
% \subsection{各领域失败分类法}

The batch runner records one terminal status per run, so the per-domain failure
\emph{categories} cannot be derived from the released logs alone. We therefore
report (i) the run-level outcome anchors per domain (Table~\ref{tab:full-domain-results})
and (ii) a triage tool that recovers the categories from each run's judge log and
agent session, which accompanies this benchmark
(\texttt{eval\_triage/05\_batch\_triage.py}, with an explicit UNKNOWN share so that
unclassified runs are never silently folded into another bucket). Applying it to
the \Neval{}-run batch yields the per-domain category matrix in
Table~\ref{tab:failure-by-domain} (PASS runs are out of scope here; rows sum to
the FAIL+ERROR count per domain). Class abbreviations: TO=ERROR\_TIMEOUT,
INFRA=ERROR\_INFRA, NOSUB=NO\_SUBMIT, FMT=FORMAT\_CONTRACT, RT=ASSERT\_RUNTIME,
SEM=ASSERT\_SEMANTIC, ANTI=ASSERT\_ANTIHARDCODE, UNK=UNKNOWN; F+E=FAIL+ERROR.
The UNKNOWN column is dominated by the ERROR-status runs (aborted pre-judging,
counted in mode (i) of \S\ref{sec:failure-taxonomy} by terminal status) whose
abort reason carries no timeout/infrastructure signature; the remaining
FAIL-with-submit UNKNOWN (49 runs, 7.0\% of the batch) is detailed in
\S\ref{sec:failure-taxonomy}. A stacked-bar view and a Cohen's $\kappa$ agreement
statistic between the tool and a human labeler on a 100-run sample are deferred
to the camera-ready version.
% 批次 runner 每个 run 只记录一个终态，因此仅凭已发布的日志无法推导各领域的失败\emph{类别}。
% 故我们给出：(i) 各领域 run 级结果锚点（表~\ref{tab:full-domain-results}）；
% (ii) 一个逐 run 归因工具，它从每个 run 的判分日志与智能体会话中恢复类别
% （eval_triage/05_batch_triage.py，并显式报告 UNKNOWN 比例，避免未分类 run 被悄悄并入别的桶）。
% 将其应用于 \Neval{}-run 批次得到表~\ref{tab:failure-by-domain} 的逐域类别矩阵（PASS run 不在此表
% 范围；各行加和为该域 FAIL+ERROR 数）。列缩写：TO=ERROR_TIMEOUT, INFRA=ERROR_INFRA,
% NOSUB=NO_SUBMIT, FMT=FORMAT_CONTRACT, RT=ASSERT_RUNTIME, SEM=ASSERT_SEMANTIC,
% ANTI=ASSERT_ANTIHARDCODE, UNK=UNKNOWN; F+E=FAIL+ERROR。UNKNOWN 列以 ERROR 终态 run（判分前中止，
% 按 \S\ref{sec:failure-taxonomy} 模式 (i) 的终态计入）为主，其中止原因无 timeout/infra 签名；
% 其余 FAIL-with-submit UNKNOWN（49 run，占批次 7.0\%）见 \S\ref{sec:failure-taxonomy}。
% 堆叠柱状图与工具对人工标注者在 100-run 样本上的 Cohen's kappa 一致性推迟到 camera-ready。

\begin{table}[ht]
\centering
\caption{Per-domain failure-category matrix for the \Neval{}-run batch, recovered by \texttt{eval\_triage/05\_batch\_triage.py} from per-run judge logs. PASS runs are excluded; each row sums to that domain's FAIL+ERROR count. UNKNOWN is reported as its own column, not folded into another class. Pass\% is PASS/(PASS+FAIL) per domain (62/620 = 10.0\% overall).}
\label{tab:failure-by-domain}
\resizebox{\columnwidth}{!}{%
\begin{tabular}{lrrrrrrrrr}
\toprule
Domain & TO & INFRA & NOSUB & FMT & RT & SEM & ANTI & UNK & F+E \\
\midrule
analog\_context & 0 & 0 & 3 & 1 & 9 & 20 & 0 & 18 & 51 \\
characterization & 0 & 0 & 5 & 0 & 7 & 8 & 0 & 3 & 23 \\
digital & 0 & 0 & 3 & 0 & 7 & 28 & 0 & 8 & 46 \\
eda\_other & 0 & 0 & 5 & 0 & 3 & 16 & 0 & 7 & 31 \\
layout\_aether & 0 & 0 & 4 & 0 & 4 & 25 & 0 & 4 & 37 \\
layout\_virtuoso & 0 & 0 & 4 & 1 & 7 & 20 & 0 & 3 & 35 \\
pg\_layout & 0 & 0 & 3 & 0 & 7 & 11 & 0 & 3 & 24 \\
physical\_verify & 0 & 0 & 8 & 0 & 18 & 37 & 0 & 7 & 70 \\
pnr\_signoff & 0 & 0 & 2 & 0 & 1 & 16 & 0 & 0 & 19 \\
schematic\_netlist & 0 & 0 & 5 & 0 & 26 & 54 & 0 & 33 & 118 \\
spice\_sim & 0 & 0 & 8 & 0 & 26 & 48 & 0 & 12 & 94 \\
verification & 0 & 0 & 5 & 0 & 11 & 11 & 0 & 8 & 35 \\
waveform & 0 & 0 & 15 & 0 & 6 & 11 & 0 & 13 & 45 \\
\midrule
\textbf{Total} & \textbf{0} & \textbf{0} & \textbf{70} & \textbf{2} & \textbf{132} & \textbf{305} & \textbf{0} & \textbf{119} & \textbf{628} \\
\bottomrule
\end{tabular}}
\end{table}

\subsection{Token Usage and Cost Analysis}
\label{sec:token-cost}
% \subsection{Token 用量与成本分析}

We can report \emph{behavioural} cost from the archived pilot sessions. For the 30
runs that carry an authoritative \texttt{run\_result.aggregateUsage} record, the
median run consumes 0.71M input and 55.6k output tokens (mean 0.86M/52.5k; totals
25.7M/1.58M), on top of a median of 26.5 tool calls (IQR 13.8--38.0) over 20.5
iterations and 1{,}047\,s of wall-clock time (IQR 552--1{,}495). The input:output
ratio of ${\sim}13{:}1$ is the signature of a long-horizon agent loop that re-sends
its context each iteration, and it is the main cost driver of flow orchestration.
A per-task token/cost table for the full \Neval{}-run batch is pending because the
batch summary records only terminal status and wall time; the per-run session logs
are retained and will be aggregated with the same parser
(\texttt{eval\_triage/trace\_anatomy.py}) in the camera-ready version.
% 我们能报告的是\emph{行为学}成本。对 30 份带有权威 run_result.aggregateUsage 记录的 run：
% 中位每 run 消耗 0.71M input / 55.6k output token（均值 0.86M / 52.5k；合计 25.7M / 1.58M），
% 同时中位发出 26.5 次工具调用（IQR 13.8--38.0）、跨 20.5 轮迭代、耗 1,047 秒墙钟（IQR 552--1,495）。
% 约 13:1 的 input:output 比是长时程智能体回路的特征——每轮迭代都要重发上下文，这也是流程编排的主要成本来源。
% 全量 \Neval{}-run 批次的逐任务 token/成本表待补：批次摘要只记录终态与墙钟，逐 run 会话日志已留存，
% 将在 camera-ready 用同一解析器（eval_triage/trace_anatomy.py）聚合给出。
